\documentclass{scrartcl}


\usepackage{scrlayer-scrpage}
\usepackage{hyperref}
\usepackage{breakurl}
\usepackage{unicode-math}

\usepackage{minted}
\usemintedstyle{trac}

\usepackage{cite}
\usepackage{url}

\usepackage{bussproofs}
\EnableBpAbbreviations

\usepackage[english]{babel}
\nonfrenchspacing


\usepackage[final,tracking=true,stretch=10,shrink=10]{microtype}
\microtypecontext{spacing=nonfrench}

\hyphenation{Isa-belle Nieuw-en-huis meta-logic multi-set multi-sets non-unit Mini-Sat}

\title{The Alethe Proof Format}
\subtitle{An Evolving Specification and Reference}
\author{Haniel Barbosa\textsuperscript{1}
\and Mathias Fleury\textsuperscript{2}
\and Pascal Fontaine\textsuperscript{3}
\and Hans-Jörg Schurr\textsuperscript{4}}
\date{}
\publishers{
	\textsuperscript{1} Universidade Federal de Minas Gerais\\
	\textsuperscript{2} Johannes Kepler University Linz\\
	\textsuperscript{3} Université de Liège\\
	\textsuperscript{4} University of Lorraine, CNRS, Inria, and LORIA, Nancy\\
}


\begin{document}

\maketitle
\tableofcontents
\clearpage

\begin{abstract}
\section*{Foreword}
This document is a speculative specification and reference of a proof
format for SMT solvers.  The format consists of a language to express
proofs and a set of proof rules.  On the one side, the language is
inspired by natural-deduction and is based on the widely used SMT-LIB
format.  The language also includes a flexible mechanism to reason
about bound variables which allows fine-grained preprocessing proofs.
%
On the other side, the rules are structured around resolution and the
introduction of theory lemmas, in the same way as CDCL($\mathcal{T}$)-based SMT
solvers.

The specification is not yet
cast in stone, but it will evolve over time. It emerged from a list
of proof rules used by the SMT solver veriT collected in a document
called ``Proofonomicon''. Following the fate presupposed by its name,
it informally circulated among researchers interested in the proofs
produced by veriT after a few months.  We now polished this document
and gave it a respectable name.

Instead of aiming for theoretical purity, our approach
is pragmatic: the specification describes the format as it is in use
right now. It will develop in parallel with practical support for the
format within SMT solvers, proof checkers, and other tools.  We believe
it is not a perfect specification that fosters the adaption of a format,
but great tooling. This document will be a guide to develop
such tools.

Nevertheless, it not only serves as a norm to ensure compatibility
between tools, it also allows us to uncover the unsatisfactory aspects
that would otherwise be hidden deep within the nooks and crannies of
solver and checker implementations.
%
Every uncovered problem presents an opportunity
to improve the format. The authors of this document overlap with the
authors of those tools and we are committed to improve the tools, the
format, and ultimately the specification together.
This document is also an invitation to other researchers to join
these efforts. To read the reference and provide feedback, or to even
implement support for Alethe into their own tools. Please
get in touch!

\bigskip
\hspace*{\fill}The authors.
\end{abstract}


This document is a reference of the
Alethe\footnote{Alethe is a genus of small birds that occur in West Africa~\cite[wp:alethe].
The name was chosen, because it
resembles the Greek word αλήθεια (alítheia) -- truth.} proof format.  Alethe is
designed to be a flexible format to represent unsatisfiability proofs
generated by {\smt} solvers.  Alethe proofs can be consumed by other systems,
such as interactive theorem provers or proof checkers.  The design
is based on natural-deduction style structure
and rules generating and operating on first-order clauses.
%
The Alethe calculus consists of two parts: the proof
language based on {\smtlib} and a collection of proof rules.
\in{Section}[sec:alethe:language] introduces the language.  First as an
abstract language, then as a concrete syntax.
%
\in{Section}[sec:alethe:rules-generic] discusses the core
concepts behind the Alethe proof rules.
%
The Appendix \about[apx:rules] presents a list of
all proof rules currently used by {\verit}.

Alethe follows a few core design principles.  First, proofs should
be easy to understand by humans to ensure working with Alethe proofs
is easy.  Second, the language of the format should directly
correspond to the language used by the solver.  Since many solvers use the
{\smtlib} language, Alethe also uses this language.
Therefore, Alethe's base logic is the many-sorted first-order logic of {\smtlib}.
Third, the format should
be uniform for all theories used by {\smt} solvers.  With the exception
of clauses for propositional reasoning, there is no dedicated syntax for
any theory.

The format is currently used by the {\smt} solver {\verit}. If requested
by the user, {\verit} outputs a proof if it can deduce that the input problem is
unsatisfiable. In proof production mode, {\verit} supports the theory of
uninterpreted functions, the theory of linear integer and real arithmetic, and
quantifiers.
%
The Alethe proofs can be reconstructed by the \tactic{smt} tactic of
the proof assistant {\isabelle}~\cite[fleury-2019,schurr-2021]~(see
\in{Section}[sec:reconstruction]).  The \tool{SMTCoq} tool can
reconstruct an older version of the format in the
proof assistant \tool{Coq}~\cite[SMTCoq].  An effort to update the
tool to the latest version of Alethe is ongoing.
%
The {\smt} solver {\cvcfive}~\cite[barbosa-2022] (the successor of
{\cvcfour}) supports Alethe experimentally as one of its multiple proof
output formats.
%
Furthermore, there is an experimental
high-performance proof checker written in Rust.\footnote{Available at
\url{https://github.com/ufmg-smite/alethe-proof-checker}.}

In addition to this reference, the proof format has been discussed in past
publications, which provide valuable background information.  The core of
the format goes back to 2011 when two publications at the PxTP workshop
outlined the fundamental ideas behind the format~\cite[besson-2011]
and proposed rules for quantifier instantiation~\cite[deharbe-2011].
%
More recently, the format has gained support for reasoning typically used for
processing, such as skolemization, substitutions, and other manipulations of
bound variables~\cite[barbosa-2019].

\end{document}
% ================== Start cutoff

\subsection[sec:alethe:language]{The Alethe Language}

This section provides an overview of the core concepts of the
\index{Alethe}Alethe language and also introduces some notation used
throughout this chapter.
The section first introduces an abstract notation to write Alethe proofs.
Then, it introduces the concrete, {\smtlib}-based syntax.  Finally,
we show how a concrete Alethe proof can be checked.

\startexample The following example shows a simple Alethe proof
expressed in the abstract notation used in this document.
It uses quantifier instantiation and resolution to show a contradiction.
The paragraphs below describe the concepts necessary to
understand the proof step by step.

\startAlethe
\bTR\bTD 1.\eTD\bTD \ctxsep \eTD\bTD $\forall x.\, (P\,x)$ \eTD\bTD $ \proofRule{assume}$ \eTD\eTR 
\bTR\bTD 2.\eTD\bTD \ctxsep \eTD\bTD $\neg (P\,a)        $ \eTD\bTD $ \proofRule{assume}$ \eTD\eTR 
\bTR\bTD 3.\eTD\bTD \ctxsep \eTD\bTD $\neg (\forall x.\, (P\,x)) \lor (P\,a)$ \eTD\bTD $\proofRule{forall_inst}\,[(x, a)]$ \eTD\eTR 
\bTR\bTD 4.\eTD\bTD \ctxsep \eTD\bTD $\neg (\forall x.\, (P\,x)), (P\,a)$ \eTD\bTD $ (\proofRule{or}\:3)$ \eTD\eTR 
\bTR\bTD 5.\eTD\bTD \ctxsep \eTD\bTD $\bot             $ \eTD\bTD $ (\proofRule{resolution}\: 1, 2, 4)$ \eTD\eTR 
\stopAlethe

\stopparagraph

\startparagraph{Many-Sorted First-Order Logic.}
Alethe builds on the {\smtlib} language.
%
This includes its many-sorted first-order logic described in \in{Section}[sec:ti-logic].
The available sorts depend on
the selected {\smtlib} theory/logic as well as on those defined by the user, but the
distinguished $\lsymb{Bool}$ sort is always available.
%
However, Alethe also extends this logic with Hilbert's choice operator $\varepsilon$.
%
The term $\varepsilon x.\, \varphi[x]$ stands for a value $v$
such that $\varphi[v]$ is true if such a value exists. Any value is
possible otherwise.  Alethe requires that $\varepsilon$ is functional
with respect to logical equivalence: if for two formulas $\varphi$, $\psi$
that contain the free variable $x$, it holds that
$(\forall x.\,\varphi ≈ \psi)$,
then $(\varepsilon x.\, \varphi)\approx(\varepsilon x.\, \psi)$ must also hold.
Note that, choice terms can only appear in Alethe proofs, not in {\smtlib} problems.

\startparagraph{Steps.}
A \index{proof}proof in the Alethe language is an indexed list of \index{step}steps.
To mimic the concrete syntax of Alethe proofs, proof steps in the
abstract notation have the form

\startAletheS
\setupTABLE[c][4][width=14em]
\bTR\bTD $i$.\eTD\bTD $c_1,\,\dots,\, c_j$ \eTD\bTD \ctxsep \eTD\bTD
$l_1,\dots ,l_k$    \eTD\bTD $(\ruleType{rule}\; p_1,\,\dots,\, p_n)\,[a_1,\,\dots,\,a_m] $\eTD\eTR
\stopAlethe

\noindent
Each step has a unique index $i \in \mathbb{I}$, where $\mathbb{I}$ is a countable
infinite set of valid indices. In the concrete syntax all {\smtlib} symbols are
valid indices, but for examples we will use natural numbers.
%
Furthermore, $l_1, \dots ,l_k$ is a clause with the literals
$l_i$. It is the conclusion of the step. If a step has the empty clause
as its conclusion (i.e., $k = 0$) we write $\bot$.  While this
muddles the water a bit with regard to steps which have the unit clause
with the unit literal $\bot$ as their conclusion, it simplifies the
notation.  We will remark on the difference if it is relevant.  The rule
name \ruleType{rule} is taken from a set of possible proof
rules (see the Appendix~\about[apx:rules]).
Furthermore, each step has a possibly empty set of premises $\{p_1,
\dots, p_n\} \subseteq \mathbb{I}$,
and a rule-dependent and possibly
empty list of arguments $[a_1, \dots, a_m]$. The list of premises
only references earlier steps, such that the proof forms a directed
acyclic graph.  If the list of premises is empty, we will drop the
parentheses around the proof rule.
The arguments $a_i$ are either terms or tuples $(x_i,
t_i)$ where $x_i$ is a variable and $t_i$ is a term. The interpretation
of the arguments is rule specific.  The list $c_1, \dots, c_j$ is
the \index{context}{\em context} of the step.  Contexts are discussed below.
Every proof ends with a step that has the empty clause as the conclusion
and an empty context.  The list of proof rules in the appendix also uses
this notation to define the proof rules.

The example above consists of five steps. Step~4 and~5 use premises.
Since step~3 introduces a tautology, it uses no premises.  However,
it uses arguments to express the substitution $[x\mapsto a]$ used to instantiate
the quantifier. Step~4 translates the disjunction into a clause.
In the example above, the contexts are all empty.
\stopparagraph

\startparagraph{Assumptions.}
An \proofRule{assume} step introduces a term as an
\index{assumption}assumption. The proof starts with a number of
\proofRule{assume} steps. Each such step corresponds to an input
assertion. Within a subproof, additional assumptions can be introduced
too. In this case, each assumption must be discharged with an appropriate
step. The rule \proofRule{subproof} can be used to do so.  In the concrete
syntax, \ruleType{assume} steps have a dedicated command \inlineAlethe{assume}
to clearly distinguish them from normal steps that use the \inlineAlethe{step}
command (see~\in{Section}[sec:alethe:syntax]).

The example above uses two assumptions which are introduced in the first
two steps.
\stopparagraph

\startparagraph{Subproofs and Lemmas.}
Alethe uses subproofs to prove lemmas and to create and manipulate the context.
%
To prove \index{lemma}lemmas, a \index{subproof}subproof can
introduce local assumptions. The \proofRule{subproof} {\em rule} discharges the
local assumptions.
From an
assumption $\varphi$ and a formula $\psi$ proved
from $\varphi$, the \proofRule{subproof} rule deduces the clause $\neg \varphi, 
\psi$ that discharges the local assumption $\varphi$.
%
A \proofRule{subproof} step cannot use a premise from a subproof nested
within the current subproof.

Subproofs are also used to manipulate the context.
As the example below shows, the abstract notation denotes subproofs by a
frame around the steps in the subproof.  In this case the subproof
concludes with a step that does not use the \proofRule{subproof} rule,
but another rule, such as the \proofRule{bind} rule.

\startexample
This example shows a refutation of the
formula $(2 + 2) ≈ 5$. The proof uses a subproof to prove the
lemma $((2 + 2) ≈ 5) \Rightarrow 4 ≈ 5$.

\startAlethe
\bTR\bTD 1.\eTD\bTD \ctxsep \eTD\bTD $(2 + 2) ≈ 5$ \eTD\bTD $\proofRule{assume}$ \eTD\eTR 
\bTR\bTD 2.\eTD\bTD[leftframe=on] \ctxsep \eTD\bTD $(2 + 2) ≈ 5 $ \eTD\bTD $\proofRule{assume}$ \eTD\eTR 
\bTR\bTD 3.\eTD\bTD[leftframe=on] \ctxsep \eTD\bTD $(2 + 2) ≈ 4 $ \eTD\bTD $\proofRule{sum_simplify}$ \eTD\eTR 
\bTR[bottomframe=on]\bTD[bottomframe=off] 4.\eTD\bTD[leftframe=on] \ctxsep \eTD\bTD $4 ≈ 5       $ \eTD\bTD $(\proofRule{trans}\ 2, 3)$ \eTD\eTR 
\bTR\bTD 5.\eTD\bTD \ctxsep \eTD\bTD $\neg((2 + 2) ≈ 5), 4 ≈ 5$ \eTD\bTD $\proofRule{subproof}$ \eTD\eTR 
\bTR\bTD 6.\eTD\bTD \ctxsep \eTD\bTD $(4 ≈ 5)≈ \bot$ \eTD\bTD $\proofRule{eq_simplify}$ \eTD\eTR 
\bTR\bTD 7.\eTD\bTD \ctxsep \eTD\bTD $\neg((4 ≈ 5)≈ \bot), \neg(4≈ 5), \bot $ \eTD\bTD $\proofRule{equiv_pos2}$ \eTD\eTR 
\bTR\bTD 8.\eTD\bTD \ctxsep \eTD\bTD $\bot $ \eTD\bTD $(\proofRule{resolution}\ 1, 5, 6, 7)$ \eTD\eTR 
\stopAlethe

\stopexample
\stopparagraph

\startparagraph{Contexts.}
A specificity of the Alethe proofs
is the step \index{context}context.
Alethe contexts are a general mechanism to write
substitutions and to change them by attaching new elements.
A context is a possible empty list $c_1,
\dots, c_l$, where each element is either a variable or a variable-term tuple
denoted $x_i\mapsto t_i$.
%
In the first case, we say that $c_i$ {\em fixes} the
variable.  The second case is a mapping.
Throughout this chapter, $\Gamma$ denotes
an arbitrary context. 

Every context $\Gamma$ induces a substitution
$\subst(\Gamma)$. If $\Gamma$ is the empty list,
$\subst(\Gamma)$ is the empty substitution, i.e, the
identity function.
The first case fixes $x_n$ and allows the context to shadow a previously defined
substitution for $x_n$:
\startformula
\subst([c_1,\dots, c_{n-1}, x_n])
\text{ is}\subst([c_1,\dots, c_{n-1}])\text{, but }x_n\text{ maps to }x_n.
\stopformula

When $\Gamma$ ends in a mapping, the substitution is extended
with this mapping:
\startformula
\subst([c_1,\dots, c_{n-1}, x_n\mapsto t_n]) =
        \subst([c_1,\dots, c_{n-1}])\circ \{x_n\mapsto t_n\}.
\stopformula

\reference[ex:alethe:substctx]{Example}The following example illustrates this idea:
\startformula\startalign
    \NC \subst([x\mapsto 7, x \mapsto g(x)]) \NC= \{x\mapsto g(7)\} \NR
    \NC \subst([x\mapsto 7, x, x \mapsto g(x)]) \NC= \{x\mapsto g(x)\} \NR
\stopalign\stopformula

Contexts are used to express proofs about the preprocessing of terms.  The
conclusions of proof rules that use contexts always have the form
\startAletheS
\setupTABLE[c][3][width=1em]
\setupTABLE[c][2][width=4em]
\setupTABLE[c][4][width=20.7em]
\setupTABLE[c][5][width=6em]
\bTR\bTD i. \eTD\bTD $\Gamma$ \eTD\bTD \ctxsep
\eTD\bTD $t ≈ u$ \eTD\bTD ($\ruleType{rule}$, …) \eTD\eTR
\stopAlethe

where the term $t$\, is the original term, and $u$ is the term after
preprocessing.  Tautologies with contexts correspond to judgments
$\vDash_T \subst(\Gamma)(t) ≈ u$.

Formally, the context can be translated to \index{abstraction+lambda}λ-abstractions and
applications.  This is discussed in the \in{Section}[sec:alethe:semantic].

\startexample[reference=ex:ti:ctx-abstract]
This example shows a proof that uses a subproof with a context to rename a bound
variable.
\startAletheS
\setupTABLE[c][4][width=17em,align=middle]
\bTR\bTD 1.\eTD \bTD \eTD
    \bTD \ctxsep \eTD\bTD $\forall x.\,(P\,x)$
    \eTD\bTD $\proofRule{assume}$ \eTD\eTR
\bTR\bTD 2.\eTD \bTD \eTD
    \bTD \ctxsep \eTD\bTD $\neg (\forall y.\,(P\,y))$
    \eTD\bTD $\proofRule{assume}$ \eTD\eTR
     \eTD\bTD \eTD\eTR
\bTR \bTD 3.\eTD\bTD[leftframe=on]  $y, x \mapsto y$ \eTD\bTD \ctxsep
     \eTD\bTD $x ≈ y$
     \eTD\bTD $\proofRule{refl}$\eTD\eTR
\bTR[bottomframe=on] \bTD[bottomframe=off] 4.\eTD\bTD[leftframe=on]  $y, x \mapsto y$ \eTD\bTD \ctxsep
     \eTD\bTD $(P\,x) ≈ (P\,y)$
     \eTD\bTD $(\proofRule{cong}\,3) $\eTD\eTR
\bTR\bTD 5.\eTD \bTD \eTD
    \bTD \ctxsep \eTD\bTD $\forall x.\,(P\,x) ≈ \forall y.\,(P\,y)$
    \eTD\bTD $\proofRule{bind}$ \eTD\eTR
     \eTD\bTD \eTD\eTR
\bTR\bTD 6.\eTD \bTD \eTD
    \bTD \ctxsep \eTD\bTD
    %\framed[frame=off, align=middle]{
	    $\neg(\forall x.\,(P\,x) ≈ \forall y.\,(P\,y)),$\\$\neg(\forall x.\,(P\,x)), (\forall y.\,(P\,y))$
    \eTD\bTD $\proofRule{equiv_pos2}$ \eTD\eTR
     \eTD\bTD \eTD\eTR
\bTR\bTD 7.\eTD \bTD \eTD
    \bTD \ctxsep \eTD\bTD $\bot$
    \eTD\bTD $(\proofRule{resolution}\,1, 2, 5, 6)$ \eTD\eTR
     \eTD\bTD \eTD\eTR
\stopAletheS
\stopexample
\stopparagraph

\startparagraph{Implicit Reordering of Equalities.}
In addition to the explicit steps, solvers might reorder equalities,
i.e., apply symmetry of the equality predicate, without generating steps.
The sole exception is the topmost equality in the conclusion of steps
with non-empty context.  The order of the arguments of this equality can
never change.
As described above, all rules that accept a non-empty context have
a conclusion of the form $t ≈ u$.  Since the context
represents a substitution applied to the left-hand side,
this equality symbol is not symmetric.

The {\smt} solver {\verit} currently applies this freedom in a restricted
form: equalities are reordered only
when the term below the equality changes during proof search. One such
case is the instantiation of universally quantified variables. If an
instantiated variable appears below an equality, then the equality might
have an arbitrary order after instantiation.
Nevertheless, consumers of Alethe must consider the possible
implicit reordering of equalities everywhere.
\stopparagraph

\subsubsection[sec:alethe:syntax]{The Syntax}

The concrete text representation of the Alethe proofs
is based on the {\smtlib} standard. \in{Figure}[fig:proof_ex] shows an
example proof as printed by {\verit} with light edits for readability.
%
The format follows the {\smtlib} standard when possible.
%
Input problems in the {\smtlib} formats are scripts.  A {\smtlib} script
is a list of commands that manipulate the {\smt} solver.  For example,
\inlineAlethe{assert} introduces an assertion, \inlineAlethe{check-sat} starts
solving, and \inlineAlethe{get-proof} instructs the {\smt} solver to print the
proof.
%
Alethe mirrors this structure. Therefore, beside the {\smtlib}
logic and term language, it also uses commands to structure the proof.
An Alethe proof is a list of commands.

\placefigure[top][fig:proof_ex]{Example proof output. Assumptions are
  introduced;   a subproof renames bound variables; the proof finishes with
  instantiation and resolution steps.}{
	  \startAletheVerb
(assume h1 (not (p a)))
(assume h2 (forall ((z1 U)) (forall ((z2 U)) (p z2))))
...
(anchor :step t9 :args ((:= z2 vr4)))
(step t9.t1 (cl (= z2 vr4)) :rule refl)
(step t9.t2 (cl (= (p z2) (p vr4)))
         :rule cong :premises (t9.t1))
(step t9 (cl (= (forall ((z2 U)) (p z2))
                (forall ((vr4 U)) (p vr4))))
         :rule bind)
...
(step t14 (cl (forall ((vr5 U)) (p vr5)))
          :rule th_resolution :premises (t11 t12 t13))
(step t15 (cl (or (not (forall ((vr5 U)) (p vr5)))
                  (p a)))
          :rule forall_inst :args ((:= vr5 a)))
(step t16 (cl (not (forall ((vr5 U)) (p vr5))) (p a))
          :rule or :premises (t15))
(step t17 (cl) :rule resolution :premises (t16 h1 t14))
	  \stopAletheVerb
}

\placefigure[top][fig:grammar]{The proof grammar.}{%
	\startformula
	\startalign[n=3,align={right,middle,left}]
   \NC \grNT{proof}          \NC\grRule\NC \grNT{proof\_command}^{*} \NR
   \NC \grNT{proof\_command} \NC\grRule\NC \textAlethe{(assume}\; \grNT{symbol}\; \grNT{proof\_term}\,\textAlethe{)} \NR
   \NC                       \NC\grOr  \NC \textAlethe{(step}\; \grNT{symbol}\; \grNT{clause}
                                        \; \textAlethe{:rule}\; \grNT{symbol} \NR
   \NC                       \NC       \NC \quad \grNT{premises\_annotation}^{?} \NR
   \NC                       \NC       \NC \quad \grNT{context\_annotation}^{?}\;\grNT{attribute}^{*}\,\textAlethe{)} \NR
   \NC                       \NC \grOr \NC \textAlethe{(anchor :step}\; \grNT{symbol}\;
                                            \NR
   \NC                       \NC       \NC \quad \grNT{args\_annotation}^{?}\;\grNT{attribute}^{*}\,\textAlethe{)} \NR
   \NC                       \NC \grOr \NC \textAlethe{(define-fun}\; \grNT{function\_def}\,\textAlethe{)} \NR
   \NC \grNT{clause}         \NC\grRule\NC \textAlethe{(cl}\; \grNT{proof\_term}^{*}\,\textAlethe{)} \NR
   \NC \grNT{proof\_term}    \NC\grRule\NC \grNT{term}\text{ extended with } \NR
   \NC                       \NC       \NC \textAlethe{(choice (}\, \grNT{sorted\_var}\,\textAlethe{)}\; \grNT{proof\_term}\,\textAlethe{)}  \NR
   \NC \grNT{premises\_annotation}  \NC\grRule\NC \textAlethe{:premises (}\; \grNT{symbol}^{+}\textAlethe{)} \NR
   \NC \grNT{args\_annotation}      \NC\grRule\NC \textAlethe{:args}\,\textAlethe{(}\,\grNT{step\_arg}^{+}\,\textAlethe{)}  \NR
   \NC \grNT{step\_arg}     \NC\grRule\NC \grNT{symbol} \grOr
                                          \textAlethe{(}\; \grNT{symbol}\; \grNT{proof\_term}\,\textAlethe{)} \NR
   \NC \grNT{context\_annotation}      \NC\grRule\NC \textAlethe{:args}\,\textAlethe{(}\,\grNT{context\_assignment}^{+}\,\textAlethe{)}  \NR
   \NC \grNT{context\_assignment}      \NC\grRule\NC \textAlethe{(}\,\grNT{sorted_var}\,\textAlethe{)}  \NR
   \NC                                 \NC \grOr \NC \textAlethe{(:=}\, \grNT{symbol}\;\grNT{proof_term}\,\textAlethe{)} \NR
  \stopalign
  \stopformula
}

Usually, an Alethe proof is associated with a concrete {\smtlib} problem
that is proved by the Alethe proof.  This can either be a concrete
problem file or, if the incremental solving commands of {\smtlib}
are used, the implicit problem constructed at the invocation of a
\inlineAlethe{get-proof} command.
In this document, we will call this {\smtlib} problem the {\em input problem}.
All symbols declared or defined in the input problem remain declared or
defined, respectively.  Furthermore, the symbolic names introduced
by the \inlineAlethe{:named} annotation also stay valid and can be used
in the script.  For the purpose of the proof rules, terms are treated
as if proxy names introduced by \inlineAlethe{:named} annotations have been
expanded and annotations have been removed.  For example, the term
\inlineAlethe{(or (! (P a) :named baz) (! baz :foo))} and
\inlineAlethe{(or (P a) (P a))} are considered to be syntactically equal.
Here \inlineAlethe{:foo} is an arbitrary {\smtlib} annotation. 

\in{Figure}[fig:grammar] shows the grammar of the proof text. It
is based on the {\smtlib} grammar, as defined in the {\smtlib}
standard~\cite[righttext={ Appendix\nbsp B}][SMTLIB].  The non-terminals
$\grNT{attribute}$,
$\grNT{function\_def}$,
$\grNT{sorted\_var}$,
$\grNT{symbol}$, and
$\grNT{term}$
are as defined in the standard. The non-terminal $\grNT{proof_term}$
corresponds to the $\grNT{term}$ non-terminal of {\smtlib}, but is
extended with the additional production for the \inlineAlethe{choice}\index{choice} binder.

Alethe proofs are a list of commands.
The \inlineAlethe{assume} command introduces a new assumption. While this
command could also be expressed using the \inlineAlethe{step} command with
a special rule, the special semantic of an assumption justifies the
presence of a dedicated command: assumptions are neither tautological nor
derived from premises.  The \inlineAlethe{step} command, on the other hand,
introduces a derived or tautological clause.  Both commands \inlineAlethe{assume}
and \inlineAlethe{step} require an index as the first argument to later
refer back to it. This index is an arbitrary {\smtlib} symbol.
It must be unique for each \inlineAlethe{assume} and \inlineAlethe{step} command.
A special restriction applies to the \inlineAlethe{assume} commands
not within a subproof,
those reference assertions in the input {\smtlib} problem.  To simplify
proof checking, the \inlineAlethe{assume} command must use the name assigned
to the asserted formula if there is any.  For example, if the input
problem contains \inlineAlethe{(assert (! (P c) :named foo))}, then
the proof must refer to this assertion (if it is needed in the proof) as
\inlineAlethe{(assume foo (P c))}.  Note that a {\smtlib} problem can
assign a name to a term at any point, not only at its first occurrence.
If a term has more than one name, any can be picked.

The second argument of \inlineAlethe{step} and \inlineAlethe{assume}
is the conclusion of the command.
For a \inlineAlethe{step}, this term is always a clause.
To express disjunctions in {\smtlib} the \inlineAlethe{or} operator is used.
This operator, however, needs at least two arguments and cannot
represent unary or empty clauses. To circumvent this, we introduce a new
\inlineAlethe{cl} operator.  It corresponds to the standard \inlineAlethe{or}
function
extended to one argument, where it is equal to the identity, and zero
arguments, where it is equal to \inlineAlethe{false}.
The \inlineAlethe{anchor} and \inlineAlethe{define-fun} commands are used for
subproofs and sharing, respectively. The \inlineAlethe{define-fun} command
corresponds exactly to the \inlineAlethe{define-fun} command of the
{\smtlib} language.

Furthermore, the syntax uses annotations as used by {\smtlib}.  The
original {\smtlib} syntax uses the non-terminal $\grNT{attribute}$.
The Alethe syntax uses some predefined annotation.  To simplify parsing,
the order in which those must be printed is strict.
The \inlineAlethe{:premises} annotation denotes the premises and is skipped
if the rule does not require premises. If the rule carries arguments, the
\inlineAlethe{:args} annotation is used to denote them.
Anchors have two annotations: \inlineAlethe{:step} provides the name of the
step that concludes the subproof and \inlineAlethe{:args} provides the context
as sorted variables and assignments.  Note that in this annotation,
the $\grNT{symbol}$ non-terminal is intended to be a variable.
After those pre-defined annotations, the solver can use additional
annotations.  This can be used for debugging, or other purposes.
A consumer of Alethe proofs {\em must} be able to parse proofs
that contain such annotations.

\startparagraph{Subproofs}
The abstract notation denotes subproofs by marking them with a vertical
line.  To map this notation to a list of commands, Alethe \index{anchor}uses the
\inlineAlethe{anchor} command.  This command indicates the start of a subproof.
A subproof is concluded by a matching \inlineAlethe{step} command.  This step must use
a \index{rule+concluding}{\em concluding rule}
(such as \proofRule{subproof}, \proofRule{bind}, and so forth).

After the \inlineAlethe{anchor} command, the proof uses \inlineAlethe{assume} commands to list
the assumptions of the subproof.  Subsequently, the subproof is a list of
\inlineAlethe{step} commands that can use prior steps in the subproofs as premises.
% the assumptions of any parent subproof
% as premises.

In the abstract notation, every step is associated with a context.  The
concrete syntax uses anchors to optimize this.
The \index{context}context is manipulated in a nested way: if a step
pops $c_1,\dots, c_n$ from a context $\Gamma$, there is an earlier step which
pushes precisely $c_1,\dots, c_n$ onto the context.
%
Since contexts can only be manipulated by push and pop, context manipulations
are nested.  The \inlineAlethe{anchor} commands push onto the context and the
corresponding \inlineAlethe{step} commands pop from the context.
To indicate these changes to the context, every anchor is associated with a list
of fixed variables and mappings.
This list can be empty.

The \inlineAlethe{:step} annotation of the anchor command is used to indicate the \inlineAlethe{step}
command that will end the subproof.  The clause of this \inlineAlethe{step}
command is the conclusion of the subproof.   While it is possible to infer the
step that concludes a subproof from the structure of the proof, the explicit
annotation simplifies reconstruction and makes the proof easier to read.
If the subproof uses a
context, the \inlineAlethe{:args} annotation of the \inlineAlethe{anchor} command
indicates the arguments added to the context for this subproof.  The
annotation provides the sort of fixed variables.

In the example proof (\in{Figure}[fig:proof_ex]) a
subproof starts at the \inlineAlethe{anchor} command.
It ends with the \proofRule{bind} steps that finishes the
proof for the renaming of the bound variable \inlineAlethe{z2}
to \inlineAlethe{vr4}.

A further restriction applies: only the conclusion of a subproof can be used
as a premise outside the subproof. Hence, a proof checking tool can
remove the steps of the subproof from memory after checking it.

\startexample[reference=ex:ti:ctx-concrete]
This example shows the proof from \in{Example}[ex:ti:ctx-abstract]
expressed as a concrete proof.

\startSMTCode
(assume h1 (forall ((x S)) (P x)))
(assume h2 (not (forall ((y S)) (P y))))
(anchor :step t5 :args ((:= x y)))
(step t3 (cl (= x y)) :rule refl)
(step t4 (cl (= (P x) (P y))) :rule cong :premises (t3))
(step t5 (cl (= (forall ((x S)) (P x)) (forall ((y S)) (P y))))
                              :rule bind)
(step t6 (cl (= (forall ((x S)) (P x)) (forall ((y S)) (P y)))
             (not (forall ((x S)) (P x)))
             (forall ((y S)) (P y)))) :rule equiv_pos2)
(step t7 (cl) :rule resolution :premises (h1 h2 t5 t6))
\stopSMTCode
\stopexample
\stopparagraph

\tikzset{
     solver/.style={draw, thin},
     system/.style={draw, thin, rounded corners},
}

% To get the colors to work correctly I have to use an odd hack inserting
% a completely transparent character.  I use negative space to then
% move back.
\definecolor[invisible][t=0.0,a=1]
\definehspace[backspace][-0.3 em]
\placefigure[top][fig:proof-states]{Abstract view of the transitions in an Alethe proof.}{
\startalignment[middle]
\starttikzpicture[node distance=2cm, auto,>=latex', thick,scale=0.8]
    \node[solver] (unsat) {\ss Unsat mode};
    \node[system, right=of unsat] (assume) {\ss Assumptions};
    \path[->] (unsat) edge[bend left] node {{\tt\tfx get-proof}} (assume);
    \path[->] (assume) edge[loop above] node[align=center] {\tfx \inlineAlethe{assume},\\\tfx\inlineAlethe{define-fun},\\\tfx  \inlineAlethe{anchor}} (assume);
    \node[system, right=of assume] (step) {\ss Steps};
    \path[->] (assume) edge[bend left] node {\tfx \inlineAlethe{step}} (step);
    \path[->]   (step) edge[loop above] node[align=center] {\tfx \color[invisible]{.}\hspace[backspace]\inlineAlethe{step},\\ \tfx\inlineAlethe{define-fun}} (step);
    \path[->] (step) edge[above, bend left] node {\tfx \color[invisible]{.}\hspace[backspace]\inlineAlethe{anchor}} (assume);
    \path[->] (step) edge[above, bend left] node {\ss \tfx Last step} (unsat);
\stoptikzpicture\stopalignment}

\startparagraph{Alethe Proof Printing States}
\in{Figure}[fig:proof-states] shows the states of an Alethe proof
abstractly.  To generate a proof, the {\smt} solver must be in the
{\em Unsat mode}, i.e., the solver determined that the problem
is unsatisfiable.  The {\smtlib} problem script then requests the proof by
invoking the \inlineAlethe{get-proof} command.  It is possible that this command
fails. For example, if proof production was not activated up front.
If there is no error, the proof is printed and printing starts with
the assumptions.  The solver prints the proof as a list of commands
according to the state.  The states ensure one constraint is maintained:
assumptions can only appear at either the beginning of a proof or right
after a subproof is started by the \inlineAlethe{anchor} command.  They cannot
be mixed with ordinary proof steps.  This simplifies
reconstruction.  Each assumption printed at the beginning of the proof
corresponds to assertions in the input problem, up to symmetry of equality.
Proof printing concludes after the last step is printed and the solver
returns to the Unsat mode and the user can issue further commands.
Usually the last step is an outermost step
(i.e., not within a subproof) that concludes the proof by deriving
the empty clause, but this is not necessary.
The solver is allowed to print some additional, useless,
steps after the proof is concluded.
\stopparagraph

\startparagraph{Sharing and Skolem Terms}
Usually, {\smt} solvers store terms internally in an efficient manner.  A term data
structure with perfect sharing ensures that every term is stored in memory
precisely once. When printing the proof, this compact storage is unfolded.
This leads to a blowup of the proof size.

Alethe can optionally use sharing\footnote{For {\verit} this can be activated
by the command-line option {\tt\minusminus{}proof-with-sharing}.} to print common
subterms only once.  This is realized using the standard naming mechanism
of {\smtlib}. A term $t$ is annotated with a name $n$ by writing
\inlineAlethe{(! }$t$\inlineAlethe{ :named }$n$\inlineAlethe{ )}
where $n$ is a symbol. After a term is annotated with a name, the name can
be used in place of the term. This is a purely syntactical replacement.
Alethe continues to use the names already used in the input problem.
Hence, terms that already have a name in the input problem can be replaced
by that name and new names introduced in the proof must not use names
already used in the input problem.

To limit the number of names, an {\smt} solver can use the following
simple approach used by {\verit}.
Before printing the proof, it iterates over all terms of the proof and
recursively descend into the terms. It marks every unmarked subterm it
visits. If a visited term is already marked, the solver assigns a new name
to this term.  If a term already has a name, it does not descend further
into the term.
%
By doing so, it ensures that only terms that appear as
child of two different parent terms get a name.  Since a named term
is replaced with its name after it first appearance, a term that only
appears as a child of one single term does not need a distinct name.
%
Thanks to the perfect
sharing representation, testing if a term is marked takes constant time
and the overall traversal takes linear time in the proof size.

To simplify reconstruction, Alethe can optionally\footnote{For {\verit} by
using the command-line option {\tt\minusminus{}proof-define-skolems}.} define
Skolem constants as functions. In this case, the proof contains a list
of \inlineAlethe{define-fun} commands that define shorthand 0-ary functions for
the \inlineAlethe{(choice }\dots\inlineAlethe{)} terms needed. Without this option,
no \inlineAlethe{define-fun} commands are issued, and the constants are expanded.
\stopparagraph

\startparagraph{Implicit Transformations}
Overall, the following aspects are treated implicitly by Alethe.
\startitemize[packed]
	\item Symmetry of equalities that are not top-most equalities in steps with
	      non-empty context.
	\item The order of literals in the clauses.
	\item The unfolding of names introduced by
	   \inlineAlethe{(! }$t$\inlineAlethe{ :named }$s$\inlineAlethe{)}.
	\item The removal of other {\smtlib} annotations of the form
	   \inlineAlethe{(! }$t$ \dots\inlineAlethe{)}.
	\item The unfolding of function symbols introduced by
	\inlineAlethe{define-fun}.\footnote{For {\verit} this is only used when the user
	introduces {\verit} to print Skolem terms as defined functions. User defined
	functions in the input problem are not supported by {\verit} in proof production
	mode.}
\stopitemize

Alethe proofs contain steps for other aspects that are commonly left implicit, such
as renaming of bound variables, and the application of substitutions.
\stopparagraph

\subsubsection[sec:alethe:semantic]{Checking Alethe Proofs}

In this section we present an abstract procedure to check if an Alethe
proof is \index{well-formed}well-formed and valid.  An Alethe proof is
well-formed only if its anchors and steps are balanced.  To check that
this is the case, we replace innermost subproofs by holes until there is
no subproof left.  If the resulting proof is free of anchors and steps
that use concluding rules, then the proof is well-formed.
If all the steps in the subproofs adhere to the conditions of
their rules, the subproof is valid.  If all subproofs are valid,
then the entire proof is valid.

Formally, an Alethe proof $P$\, is a list $[C_1, \dots, C_n]$ of steps
and anchors.
Since every step uses an unique index, we assume that each step
$C_i$ in $P$\, uses $i$ as its index.
The context
only changes at anchors and subproof-concluding steps.
Therefore, the elements of $C_1, \dots, C_n$ that are steps
are not associated with a context.  
Instead, the context can be computed
from the prior anchors.
The anchors only ever extends the context. 

To check an Alethe proof we can iteratively eliminate the first-innermost
subproof, i.e., the innermost subproof that does not come after a
complete subproof.  The restriction to the first subproofs simplifies
the calculation of the context of the steps in the subproof.

\startdefinition{First-Innermost Subproof}
  Let $P$\, be the proof $[C_1, \dots, C_n]$ and $1 \leq \mathit{start}
  < \mathit{end} \leq n$ be two indices such that
  \startitemize[packed]
    \item $C_{\mathit{start}}$ is an anchor,
    \item $C_{\mathit{end}}$ is a step that uses a concluding rule,
    \item no $C_k$ with $k < \mathit{start}$ uses a concluding rule,
    \item no $C_k$ with $\mathit{start} < k < \mathit{end}$ is an anchor or
     a step that uses a concluding rule.
  \stopitemize
  Then $[C_{\mathit{start}}, \dots, C_{\mathit{end}}]$ is the
  \index{subproof+first-innermost}first-innermost subproof of $P$.
\stopdefinition

\startexample
The proof in \in{Example}[ex:ti:ctx-concrete] has only one subproof
and this subproof is also a first-innermost subproof.  It is the subproof
\startSMTCode
(anchor :step t5 :args ((:= x y)))
(step t3 (cl (= x y)) :rule refl)
(step t4 (cl (= (P x) (P y))) :rule cong :premises (t3))
(step t5 (cl (= (forall ((x S)) (P x)) (forall ((y S)) (P y))))
                              :rule bind)
\stopSMTCode
\stopexample

It is easy to calculate the context of the first-innermost subproof.

\startdefinition{Calculated Context}
  Let $[C_{\mathit{start}}, \dots, C_{\mathit{end}}]$ be
  the first-innermost subproof of $P$.
	For $\mathit{start} \leq i < \mathit{end}$, let
  $A_1, \dots, A_m$ be the anchors among $C_1, \dots, C_{i-1}$.

  The \index{context+calculated}calculated context of $C_i$ is the context
  \startformula
  c_{1,1}, \dots, c_{1, n_1}, \dots, c_{m,1}, \dots, c_{m, n_m}
  \stopformula
  where $c_{k,1}, \dots, c_{k, n_k}$ is the list of fixed variables
  and mappings associated with $A_k$.
\stopdefinition

Note that if $C_i$ is an anchor, its calculated context does not
contain the elements associated with $C_i$.
Therefore, the context of $C_{\mathit{start}}$
is the context of the steps before the subproof.
Furthermore, the step $C_{\mathit{end}}$ is the concluding
step of the subproof and must have the same context as the steps surrounding
the subproof.
Hence, the context of $C_{\mathit{end}}$ is the calculated
context of $C_{\mathit{start}}$.

\startexample
The calculated context of the steps \color[SmtStepId]{\tt t3} and \color[SmtStepId]{\tt t5} in
\in{Example}[ex:ti:ctx-concrete] is the context $x \mapsto y$.
The calculated context of the concluding step \color[SmtStepId]{\tt t5} and the anchor is empty.
\stopexample

\noindent
A first-innermost subproof is valid if all its steps adhere to
the conditions of their rule and only use premises that occur in front of them
in the subproof.  The conditions of each rule are listed the
Appendix~\about[apx:rules].

\startdefinition{Valid First-Innermost Subproof}
	Let $[C_{\mathit{start}}, \dots, C_{\mathit{end}}]$
	be the first-innermost subproof of $P$.
	The subproof is \index{subproof+valid}{\em valid} if
  \startitemize[packed]
    \item all steps $C_i$ with $\mathit{start} < i <
    \mathit{end}$ only use premises $C_j$ with $\mathit{start} <
    j < i$,
    \item all $C_i$ that are steps adhere to the conditions of their
    rule under the calculated context of $C_i$,
    \item the step $C_{\mathit{end}}$
		adheres to the conditions of its
    rule under the calculated context of $C_{\mathit{start}}$.
  \stopitemize
\stopdefinition

Since the \proofRule{assume} rule expects an empty context, an admissible
subproof can contain assumptions only if the rule used by $C_{\mathit{end}}$
is the \proofRule{subproof} rule.

To eliminate a subproof we can replace the subproof with a hole that has
at its conclusion the conclusion of the subproof.  This is safe as long
as the subproof that is eliminated is valid (see \in{Section}[sec:alethe:soundness-eh]).

\startdefinition
The function $E$ eliminates the first-innermost subproof from a proof
if there is one.
Let $P$\, be a proof $[C_1, \dots C_n]$.
Then $E(P) = P$\, if $P$\, has no first-innermost subproof.
Otherwise, $P$\, has the first-innermost subproof
$[C_{\mathit{start}}, \dots, C_{\mathit{end}}]$, and
$E(P) = [C_1, \dots, C_{\mathit{start}-1}, C', C_{\mathit{end}+1},
\dots, C_n]$ where $C'$ is a new step that uses the \proofRule{hole} rule
and has the index, conclusion, and premises of $C_{\mathit{end}}$.
\stopdefinition

It is important to add the premises of $C_{\mathit{end}}$
to $C'$.  The \proofRule{let} rule can use additional premises
and omitting those premises results in an unsound step.
We can apply $E$\, iteratively to a proof $P$\, until we reach the least fixed point.
Since $P$ is finite we will always reach a fixed point in finitely many steps.
The result is a list $[P_0, P_1, P_2, \dots, P_{\mathit{last}}]$ where $P_0 = P$,
$P_1 = E(P)$, $P_2 = E(E(P))$ and $P_{\mathit{last}} = E(P_{\mathit{last}})$.

\startexample
Applying $E$\, to the proof in
\in{Example}[ex:ti:ctx-concrete] gives us the proof

\startSMTCode
(assume h1 (forall ((x S)) (P x)))
(assume h2 (not (forall ((y S)) (P y))))
(step t5 (cl (= (forall ((x S)) (P x)) (forall ((y S)) (P y))))
                                      :rule hole)
(step t6 (cl (= (forall ((x S)) (P x)) (forall ((y S)) (P y)))
             (not (forall ((x S)) (P x)))
             (forall ((y S)) (P y)))) :rule equiv_pos2)
(step t7 (cl) :rule resolution :premises (h1 h2 t5 t6))
\stopSMTCode

Since this proof contains no subproof, it is also $P_{\mathit{last}}$.
\stopexample

\startdefinition[reference=def:well_formed_proof,title={Well-Formed Proof}]
  The Alethe proof $P$\, is \index{proof+well-formed}well-formed
  if every step uses a unique index and $P_{\mathit{last}}$
  contains no anchor or step that uses a concluding rule.
\stopdefinition

\startdefinition{Valid Alethe Proof}
	The proof $P$\, is a \index{proof+valid}{\em valid Alethe proof} if
  \startitemize[packed]
  	\item $P$\, is well-formed,
  	\item $P$\, does not contain any step that uses the \proofRule{hole} rule,
    \item $P_{\mathit{last}}$ contains a step that has the empty clause as its conclusion,
		\item the first-innermost subproof of every $P_i$, $i < \mathit{last}$ is valid,
    \item all steps $C_i$ in $P_{\mathit{last}}$ only use premises
          $C_j$ in $P_{\mathit{last}}$ with $1 \leq j < i$,
    \item all steps $C_i$ in $P_{\mathit{last}}$ adhere to the conditions of their
    rule under the empty context.
  \stopitemize
\stopdefinition

The condition that $P$\, contains no hole ensures that the original
proof is complete and holes are only introduced by eliminating valid
subproofs.

\startexample
The proof in \in{Example}[ex:ti:ctx-concrete] is valid.  The only
subproof is valid, the proof contains no hole, and $P_{\mathit{last}}$
contains the step {\ss t7} that concludes with the empty clause.
\stopexample

It is sometimes useful to speak about the steps that are not within a
subproof.  We call such a step an \index{step+outermost}{\em outermost
step}.  In a well-formed proof those are the steps
of $P_{\mathit{last}}$.

\subsubsection{Contexts and Metaterms}

We now direct our attention to subproofs with \index{context}contexts.
It is useful to give a precise semantic
to contexts to have the tools to check that rules that use contexts
are sound.
Contexts are local in the sense that they affect only the
proof step they are applied to. 
For the full details on contexts see~\cite[barbosa-2019].
The presentation here is adapted from this publication, but omits some
details.

To handle subproofs with contexts, we translate the contexts into
λ-terms.
%
This allows us to
leverage the \index{lambda calculus}λ-calculus as an existing well-understood theory of \index{binder}binders.
%
These λ-terms\index{term+lambda} are called \index{term+meta}\index{metaterms}{\em metaterms}.

\startdefinition{Metaterm}
  Metaterms are expressions constructed by the grammar
  \startformula
    M \,::=\, \groundbox{$t$}\,\mid\, \lambda x.\,M\,\mid\,(\lambda \bar{x}_n.\,
    M)\,\bar{t}_n
  \stopformula
where $t$ is an ordinary term and $t_i$ and $x_i$ have matching sorts for
all $0 \leq i \leq 1$.
\stopdefinition

According to this definition, a metaterm is either a boxed term, a
\index{abstraction+lambda}λ-abstraction, or an application to an uncurried λ-term.
The annotation $\groundbox{$t$}$ delimits terms from the context, a simple
λ-abstraction is used to express fixed variables, and the
application expresses simulations substitution of $n$ terms.\footnote{
	The box annotation used here is unrelated to the boxes
	within the {\smt} solver discussed in the introduction.}

We use $=_{\alpha\beta}$ to denote syntactic equivalence modulo
α-equivalence and β-reduction.

Proof steps with contexts can be encoded into proof steps with empty
contexts, but with metaterms.  A proof step

\startAletheS
\setupTABLE[c][2][width=2em]
\setupTABLE[c][3][width=1em]
\setupTABLE[c][4][width=22.7em]
\setupTABLE[c][5][width=6em]
\bTR\bTD i. \eTD\bTD $\Gamma$ \eTD\bTD \ctxsep
\eTD\bTD $t ≈ u$ \eTD\bTD $(\ruleType{rule}\; \bar{p}_n)\;[\bar{a}_m]$ \eTD\eTR
\stopAlethe

\noindent
is encoded into

\startAletheS
\setupTABLE[c][2][width=2em]
\setupTABLE[c][3][width=1em]
\setupTABLE[c][4][width=22.7em]
\setupTABLE[c][5][width=6em]
\bTR\bTD i. \eTD\bTD  \eTD\bTD \ctxsep
\eTD\bTD $L(\Gamma)[t] \simeq R(\Gamma)[u]$ \eTD\bTD $(\ruleType{rule}\; \bar{p}_n)\;[\bar{a}_m]$ \eTD\eTR
\stopAlethe

\noindent
where
\startformula
\startalign[n=2,m=2,align={right,left}, distance=1.8em]
 \NC    L(\emptyset)[t] \NC= \groundbox{$t$}                  \NC    R(\emptyset)[u] \NC= \groundbox{$u$} \NR
 \NC L(x, \bar{c}_m)[t] \NC= \lambda x.\,L(\bar{c}_m)[t]    \NC R(x, \bar{c}_m)[u] \NC= \lambda x.\,R(\bar{c}_m)[u] \NR
  \NC L(\bar{x}_n\mapsto\bar{s}_n, \bar{c}_m)[t] \NC= (\lambda \bar{x}_n.\,L(\bar{c}_m)[t]) \bar{s}_n
     \NC R(\bar{x}_n\mapsto\bar{s}_n, \bar{c}_m)[u] \NC= R(\bar{c}_m)[u] \NR
\stopalign
\stopformula

To achieve the same effect as using the $\subst()$ function described above, we
can translate the terms into metaterms, perform β-reduction, and rename
bound variables if necessary~\cite[righttext={ Lemma~11}][barbosa-2019].

\startexample
The example on page~\ref[page][ex:alethe:substctx] becomes
\startformula
\startalign[n=1,align=left]
\NC L(x\mapsto 7, x \mapsto g(x))[x] = (\lambda x.\,(\lambda x.\,\groundbox{$x$})\,(g(x)))\, 7
=_{\alpha\beta} \groundbox{$g(7)$} \NR
\NC L(x\mapsto 7, x, x \mapsto g(x))[x] = (\lambda x.\,\lambda x.\,(\lambda x.\,\groundbox{$x$})\,(g(x)))\, 7
=_{\alpha\beta} \lambda x.\,\groundbox{$g(x)$} \NR
\stopalign
\stopformula
\stopexample

Most proof rules that operate with contexts can easily be translated into
proof rules using metaterms.  The exception are the tautologous rules,
such as \proofRule{refl} and the \ruleType{\cdots{}_simplify} rules.

Steps that use such rules always encode a judgment
$\vDash \Gamma\,\scale[grid=halfline]{▷}\,t ≈ u$.  With the encoding described above
we get $L(\Gamma)[t] ≈ R(\Gamma)[u]
=_{\alpha\beta} \lambda \bar{x}_n.\,\groundbox{$t'$} ≈
\lambda \bar{x}_n.\,\groundbox{$u'$}$ with some terms $t'$, $u'$.
To handle those terms, we use the $\reify()$ function.
This function is defined as
\startformula
\reify(\lambda \bar{x}_n.\,\groundbox{$t$} ≈
\lambda \bar{x}_n.\,\groundbox{$u$}) = \forall \bar{x}_n.\,(t ≈ u).
\stopformula
Therefore,
all tautological rules with contexts represent a judgment\\
$\vDash \reify(T ≈ U)$
where $T =_{\alpha\beta} L(\Gamma)[t]$ and $U =_{\alpha\beta} R(\Gamma)[u]$.

\startexample
  Consider the step

\startAletheS
\setupTABLE[c][3][width=1em]
\setupTABLE[c][2][width=4em]
\setupTABLE[c][4][width=20.7em]
\setupTABLE[c][5][width=6em]
\bTR\bTD i. \eTD\bTD $y, x \mapsto y$ \eTD\bTD \ctxsep
\eTD\bTD $x + 0 ≈ y$ \eTD\bTD $\proofRule{sum_simplify}$ \eTD\eTR
\stopAlethe

\noindent
Translating the context into metaterms leads to

\startAletheS
\setupTABLE[c][3][width=1em]
\setupTABLE[c][2][width=4em]
\setupTABLE[c][3][width=1em]
\setupTABLE[c][4][width=20.7em]
\setupTABLE[c][5][width=6em]
\bTR\bTD i. \eTD\bTD \eTD\bTD \ctxsep
\eTD\bTD $(\lambda y.\,(\lambda x.\,\groundbox{$x + 0$})\, y) ≈
        (\lambda y.\,\groundbox{$y$})$ \eTD\bTD $\proofRule{sum_simplify}$ \eTD\eTR
\stopAlethe

\noindent
Applying β-reduction leads to

\startAletheS
\setupTABLE[c][3][width=1em]
\setupTABLE[c][2][width=4em]
\setupTABLE[c][3][width=1em]
\setupTABLE[c][4][width=20.7em]
\setupTABLE[c][5][width=6em]
\bTR\bTD i. \eTD\bTD \eTD\bTD \ctxsep
\eTD\bTD $(\lambda y.\,\groundbox{$y + 0$}) ≈
        (\lambda y.\,\groundbox{$y$}) $ \eTD\bTD $\proofRule{sum_simplify}$ \eTD\eTR
\stopAlethe

\noindent
Finally, using $\reify()$ leads to

\startAletheS
\setupTABLE[c][3][width=1em]
\setupTABLE[c][2][width=4em]
\setupTABLE[c][3][width=1em]
\setupTABLE[c][4][width=20.7em]
\setupTABLE[c][5][width=6em]
\bTR\bTD i. \eTD\bTD \eTD\bTD \ctxsep
\eTD\bTD $\forall y.\,(y + 0 ≈  y)$ \eTD\bTD $\proofRule{sum_simplify}$ \eTD\eTR
\stopAlethe

\noindent
This obviously holds in the theory of linear arithmetic.
\stopexample


\subsubsection[sec:alethe:soundness-eh]{Soundness}

Any proof calculus should be sound.  In the case of Alethe, most proof
rules are
standard rules, or simple tautologies.  The rules that use context
are unusual, but those proof rules were previously shown to be
sound~\cite[barbosa-2019].
Alethe does not use any rules that are merely satisfiability preserving.
The skolemization rules replace the bound variables with choice terms
instead of fresh symbols.\footnote{The \inlineAlethe{define-fun} function
can introduce fresh symbols, but we will assume here that those
commands have been eliminated by unfolding the definition.}
All Alethe rules express semantic implications.
Overall, we assume in this document that the proof rules and proofs
written in the abstract notation are sound.

Nevertheless, a modest gap remains.  The concrete, command-based
syntax does not precisely correspond to the abstract notation.
In this section we will address the soundness of concrete Alethe
proofs.

\starttheorem[reference=thm:sound,title={Soundness of Concrete Alethe Proofs}]
	If there is a valid Alethe proof $P = [C_1, \dots, C_n]$ that has the formulas
	$\varphi_1, \dots, \varphi_m$ as the conclusions of the outermost \proofRule{assume}
	steps, then
	\startformula
	\varphi_1, \dots, \varphi_m \vDash \bot.
	\stopformula
\stoptheorem

Here, $\vDash$ represents
semantic consequence in the many-sorted first order logic of {\smtlib}
with the theories of uninterpreted functions and linear arithmetic extended
to clauses.

To show the soundness of a valid Alethe proof $P = [C_1, \dots, C_n]$,
we can use the same approach as for the definition of validity: consider
first-innermost subproof first and then replace them by holes.
%
Since valid proofs do not contain holes, we have to generalize the induction
to allow holes that were introduced by the elimination of subproofs.
%
We start with simple subproofs with empty contexts and without
nested subproofs.

\startlemma[reference=lem:sound_subproof]
  Let $P$\, be a proof that contains a valid first-innermost subproof where
  $C_{\mathit{end}}$ is a \proofRule{subproof} step.  Let
  $\psi$ be the conclusion of $C_{\mathit{end}}$.
  Then $\vDash \psi$ holds.
\stoplemma
\startproof
  First, we use induction on the number of steps $n$ after
  $C_{\mathit{start}}$.  Let $\psi_n$ be the conclusion
  of $C_{\mathit{start}+n}$ and $V_n$ the conclusions of
  the \proofRule{assume} steps in $[C_{\mathit{start}}, \dots,
  C_{\mathit{start}+n}]$.
  Assumptions always introduce unit clauses.  We will identify
  unit clauses with their single literal.
  We will
  show $V_n \vDash \psi_n$ if $\mathit{start}+n < \mathit{end}$.

  If $n = 1$, then $C_{\mathit{start}+n} = C_{\mathit{start}+1}$ must
  either be a tautology, or an \proofRule{assume} step.  In the first case,
  $\vDash \psi_{\mathit{start}+1}$ holds, and in the second case
  $\psi_{\mathit{start}+1} \vDash \psi_{\mathit{start}+1}$ holds.

  For subsequent $n$, $C_{\mathit{start}+n}$ is
  either an ordinary step, or an \proofRule{assume} step.  In the second case,
  $\psi_{\mathit{start}+n} \vDash \psi_{\mathit{start}+n}$ which can
  be weakened to $V_n \vDash \psi_{\mathit{start}+n}$.
  In the first case, the step $C_{\mathit{start}+n}$ has a set of
  premises $S$.  
  For each step $C_{\mathit{start}+i} \in S$ we have $i < n$ and
  $V_i \vDash \psi_{\mathit{start}+i}$ due to the induction
  hypothesis. Since $i < n$, the premises $V_i$ are a subset of $V_n$ and
  we can weaken $V_i \vDash \psi_{\mathit{start}+i}$
  to $V_n \vDash \psi_{\mathit{start}+i}$.  Since all premises of
  $C_{\mathit{start}+n}$ are the consequence of $V_n$ we get
  $V_n \vDash \psi_n$.

  The step $C_{\mathit{end}-1}$ is the last step of the subproof that
  does not use a concluding rule.
	At this step we have $V_{\mathit{end}-1} \vDash \psi_{\mathit{end}-1}$.
	Since $C_{\mathit{end}}$ is not an \proofRule{assume} step, the
	set $V_{\mathit{end}-1} = \{\varphi_1, \dots, \varphi_m\}$ contains
	all assumptions in the subproof.
	By the deduction theorem we get
	\startformula
	\vDash \varphi_1 \land \cdots \land \varphi_m \limp \psi_{\mathit{end}-1}.
	\stopformula
	This can be transformed into the clause
	\startformula
	\vDash \neg\varphi_1, \cdots, \neg\varphi_m, l_1, \dots, l_o.
	\stopformula
	where $l_1, \dots, l_o$ are the literals of $\psi_{\mathit{end}-1}$.
	This clause is exactly the conclusion of the 
	step $C_{\mathit{end}}$
	according to the definition of the \proofRule{subproof} rule.
\stopproof

\noindent
We can do the same reasoning as for \in{Lemma}[lem:sound_subproof] for
subproofs with contexts.  This is slightly complicated by the
\proofRule{let} rule that can use extra premises.

\startlemma[reference=lem:sound_subproof_context]
  Let $P$\, be a proof that contains a
  valid first-innermost subproof where
  $C_{\mathit{end}}$ is a step using one of:
  \proofRule{bind}, \proofRule{sko_ex}, \proofRule{sko_forall},
  \proofRule{onepoint}, \proofRule{let}.

  Then $V \vDash \Gamma \vartriangleright \psi$ where $\Gamma$ is the
  calculated context of $C_{\mathit{start}}$ and $\psi$ is the conclusion
  of $C_{\mathit{end}}$.  The set $V$\, is empty if $C_{\mathit{end}}$
  does not use the \proofRule{let} rule.  Otherwise, it contains all
  conclusions of the
  \proofRule{assume} steps among $[C_{δ}, \dots,
  C_{\mathit{start}}]$ where $δ$ is either the largest index
  $δ < {\mathit{start}}$ such
  that $s_{δ}$ is an anchor, or $1$ if no such index exist.
  Hence, there is no anchor between $C_{δ}$ and $C_{\mathit{start}}$. 
\stoplemma
\startproof
  The step $C_{\mathit{start}}$ is an anchor due to the definition
  of innermost-first subproof.
  Let $c_1, \dots, c_n$ be the context introduced by the anchor
  $C_{\mathit{start}}$, and let $\Gamma$ be the calculated context
  of $C_{\mathit{start}}$.
  $\Gamma' = \Gamma, c_1, \dots, c_n$.
	is the calculated context of the steps in the subproof after
  $C_{\mathit{start}}$.

  The step $C_{\mathit{end}}$ is a step

\startAletheS
\setupTABLE[c][1][width=4em]
\setupTABLE[c][4][width=19.0em]
\setupTABLE[c][5][width=6em]
\bTR\bTD \eTD\bTD[leftframe=on] \eTD\bTD 
\eTD\bTD \cdots \eTD\bTD \eTD\eTR
\bTR[bottomframe=on]\bTD[bottomframe=off] $\mathit{end}-1$. \eTD\bTD[leftframe=on] $\Gamma'$ \eTD\bTD \ctxsep
\eTD\bTD $\,\psi'$ \eTD\bTD $(\dots)$ \eTD\eTR
\bTR\bTD $\mathit{end}$. \eTD\bTD $\Gamma$ \eTD\bTD \ctxsep
\eTD\bTD $\psi$ \eTD\bTD $(\ruleType{rule}\;i_1, \dots, i_n)$ \eTD\eTR
\stopAlethe

  Since we assume the step $C_{\mathit{end}}$ is correctly employed,
  $\vDash \Gamma \vartriangleright \psi$ holds, as long as
  $\vDash \Gamma' \vartriangleright \psi'$ holds.

  We perform the same induction as for
  \in{Lemma}[lem:sound_subproof] over the steps in
  $[C_{\mathit{start}}, \dots, C_{\mathit{end}}]$.
  %
  Since $C_{\mathit{end}}$ does not use the \proofRule{subproof} rule,
  the subproof does not contain any assumptions and
  $V_i$ stays empty.
  %
  Again, we are interested in the step $C_{\mathit{end}-1}$.  At this step we
  get $\vDash \Gamma' \vartriangleright \psi'$.

  Only the \proofRule{let} rule uses additional premises
  $C_{i_1}, \dots, C_{i_n}$.  Hence, for all other rules, the conclusion
  cannot depend on any step outside the subproof and $V$\, is empty.
  % TODO: this is ugly, let is ugly
  Due to the definition of first-innermost subproof, all steps $C_{i_1},
  \dots, C_{i_n}$ are in the same subproof that starts at $C_{δ}$.

  The steps $C_{i_1}, \dots, C_{i_n}$ might
  depend on some \proofRule{assume} steps that appear before them
  in their subproof.  This is the case if the steps are outermost steps, or
  if the subproof that starts at $C_{δ}$ concludes with
  a \proofRule{subproof} step. 
  %
  In this case we can, as we saw in the proof of
  \in{Lemma}[lem:sound_subproof], weaken their judgments to include
  all assumptions in $[C_{δ}, \dots, C_{\mathit{start}}]$.

  If the subproof that starts at $C_{δ}$ concludes with
  any other rule, then there cannot be any assumptions and $V$\, is empty.
\stopproof

By using \in{Lemma}[lem:sound_subproof] and
\in{Lemma}[lem:sound_subproof_context] we can now show that
a valid, concrete Alethe proof is sound. That is, we can show
\in{Theorem}[thm:sound].

\startproof
	Since $P =[C_1, \dots, C_n]$ is valid, all steps that do not use the
	\proofRule{hole} rule adhere to their rule.  Since we assume that the
	abstract notation and the rules are sound, we only have to
	worry about the steps using the \proofRule{hole} rule.
	Those should be sound, i.e., for a \proofRule{hole} step with the conclusion
	$\psi$, premises $V$, and context $\Gamma$
	the judgment $V \vDash \Gamma \vartriangleright \psi$ must hold.

	Since $P$\, is a valid proof there is a sequence
	$[P_0, \dots, P_{\mathit{last}}]$ as discussed in \in{Section}[sec:alethe:semantic].
	For $i < \mathit{last}$, $E(P_i) = P_{i+1}$ replaces the
	first-innermost subproof in $P_i$ by a hole with the conclusion
	$\psi$.  Furthermore, the context of the introduced hole
	corresponds to the context $\Gamma$ of the start of the subproof.
	Since $P$ is a valid proof, the first-innermost subproof eliminated by $E$ is
	always valid.
	Therefore,
	we can apply \in{Lemma}[lem:sound_subproof]
	or \in{Lemma}[lem:sound_subproof_context] to conclude that the hole introduced
	by $E$ is sound.

	Since $P_0$ does not contain any holes, the holes in each proof
	$P_i$ are all introduced by innermost-first subproof elimination.
	Therefore, they are sound. In consequence, all holes in $P_{\mathit{last}}$ are
	sound and we can perform the same
	argument as
	in the proof of \in{Lemma}[lem:sound_subproof] to the proof
	$P_{\mathit{last}}$.

	Let $j$ be the index of the step in $P_{\mathit{last}}$ that concludes
	with the empty clause.
	Let $\mathit{start} = 1$
	and $\mathit{end} = j$ in the argument of \in{Lemma}[lem:sound_subproof].
	This shows that $V \vDash \bot$, where $V$\, is the
	conclusion of the \proofRule{assume} steps in the sublist $[C_1, \dots, C_j]$
	of $P_{\mathit{last}}$.  We can weaken
	this by adding the conclusions of the \proofRule{assume} steps in
	$[C_j, \dots, C_n]$ of $P_{\mathit{last}}$
	to get $\varphi_1, \dots, \varphi_m \vDash \bot$.
\stopproof

\subsection[sec:alethe:rules-generic]{The Alethe Rules}

Together with the language, the Alethe format also includes a set
of proof rule.
The Appendix~\about[apx:rules]
gives a full list of all
proof rules.  Currently, the proof rules correspond to the rules
that the solver {\verit} can emit.  For the rest of this section, we will discuss some
general concepts related to the rules.

\startparagraph{Tautologous Rules and Simple Deduction.}
Most rules introduce tautologies. One example is
the \proofRule{and_pos} rule: $\neg (\varphi_1 \land \varphi_2 \land
\dots \land \varphi_n), \varphi_i$.
%
Other rules derive their conclusion from a single premise.
%
Those rules are primarily used to simplify Boolean
connectives during preprocessing. For example, the \proofRule{implies}
rule eliminates an implication: From $\varphi_1 \limp \varphi_2$,
it deduces $\neg \varphi_1, \varphi_2$.
\stopparagraph

\startparagraph{Resolution.}
{\cdclt}-based {\smt} solvers, and especially their {\sat} solvers,
are fundamentally based on resolution of clauses.
Hence, Alethe also has dedicated clauses and a resolution proof
rule.  However, since {\smt} solvers do not enforce a strict
clausal normal form, ordinary disjunction is also used.
Keeping clauses and disjunctions distinct simplifies rule checking.
For example, in the case of
resolution there is a clear distinction between unit clauses where
the sole literal starts with a disjunction and non-unit clauses.  The
syntax for clauses uses the \inlineAlethe{cl} operator, while disjunctions
use the standard {\smtlib} \inlineAlethe{or} operator. The \proofRule{or}
\emph{rule} is responsible for converting disjunctions into clauses.

The Alethe proofs use a generalized propositional
resolution rule with the name \proofRule{resolution} or
\proofRule{th_resolution}. Both names denote the same rule. The
difference only serves to distinguish if the rule was introduced by
the {\sat} solver or by a theory solver. The resolution step is purely
propositional; there is no notion of a unifier.  The resolution
rules also implicitly simplifies repeated negations at the head
of literals.

The premises of a resolution step are clauses, and the conclusion
is a clause that has been derived from the premises by some binary
resolution steps.
\stopparagraph

\startparagraph[reference=rule:forall_inst,title={Quantifier Instantiation.}]
To express quantifier instantiation, the rule \proofRule{forall_inst}
is used. It produces a formula of the form $(\neg \forall \bar
x_n.\,\varphi) \lor \varphi[\bar t_n]$, where $\varphi$ is
a term containing the free variables $\bar x_n$, and for each $i$ the
ground term $t_i$ is a new term with the same sort as $x_i$.\footnote{
	For historical reasons, \proofRule{forall} has a unit clause with a disjunction
	as its conclusion and not the clause $(\neg \forall \bar x_n.\,\varphi), \varphi[\bar t_n]$.
}

The arguments of a \proofRule{forall_inst} step is the list $(x_1 , t_1),
\dots, (x_n, t_n)$. While this information can be recovered from the term,
providing it explicitly helps reconstruction because the implicit reordering of
equalities obscures which terms have been used as instances.
Existential quantifiers are handled by skolemization.
\stopparagraph

\startparagraph{Linear Arithmetic.}
Proofs for linear arithmetic use a number of straightforward rules,
such as \proofRule{la_totality} ($t_1 \leq t_2 \lor t_2 \leq t_1$)\footnote{
	This rule also has a unit clause with a disjunction as its conclusion.}
and the main rule \proofRule{la_generic}.  The conclusion of an
\proofRule{la_generic} step is a tautology $\neg \varphi_1, \neg
\varphi_2, \dots, \neg \varphi_n$ where the $\varphi_i$ are linear
(in)equalities. Checking the validity of this clause amounts to
checking the unsatisfiability of the system of linear equations
$\varphi_1, \varphi_2, \dots, \varphi_n$.  The annotation of an
\proofRule{la_generic} step contains a coefficient for each (in)equality.
The result of forming the linear combination of the literals with
the coefficients is a trivial inequality between constants.

\startexample
The following example is the proof for the unsatisfiability
of $(x+y<1) \lor$ $(3<x)$, $x≈ 2$, and $0≈ y$.

\startAlethe
\setupTABLE[c][3][width=19em,align=middle]
\bTR\bTD 1.\eTD\bTD \ctxsep \eTD\bTD $(3 < x) \lor (x + y < 1) $ \eTD\bTD      $ \proofRule{assume}$ \eTD\eTR
\bTR\bTD 2.\eTD\bTD \ctxsep \eTD\bTD $x≈ 2                   $ \eTD\bTD      $ \proofRule{assume}$ \eTD\eTR
\bTR\bTD 3.\eTD\bTD \ctxsep \eTD\bTD $0≈ y                   $ \eTD\bTD      $ \proofRule{assume}$ \eTD\eTR
\bTR\bTD 4.\eTD\bTD \ctxsep \eTD\bTD $(3 < x), (x + y < 1)        $ \eTD\bTD      $ (\proofRule{or}\,1)$ \eTD\eTR
\bTR\bTD 5.\eTD\bTD \ctxsep \eTD\bTD $\neg (3<x), \neg (x≈ 2)$ \eTD\bTD      $ \proofRule{la_generic}\, [1.0, 1.0]$ \eTD\eTR
\bTR\bTD 6.\eTD\bTD \ctxsep \eTD\bTD $\neg (3<x)                  $ \eTD\bTD      $ (\proofRule{resolution}\, 2, 5)$ \eTD\eTR
\bTR\bTD 7.\eTD\bTD \ctxsep \eTD\bTD $x + y < 1                   $ \eTD\bTD      $ (\proofRule{resolution}\, 4, 6)$ \eTD\eTR
\bTR\bTD 8.\eTD\bTD \ctxsep \eTD\bTD $\neg (x + y < 1), \neg (x≈ 2) \lor \neg (0 ≈ y)$ \eTD\bTD $\proofRule{la_generic}\,[1.0, -1.0, 1.0]$ \eTD\eTR
\bTR\bTD 9.\eTD\bTD \ctxsep \eTD\bTD $\bot      $ \eTD\bTD      $ (\proofRule{resolution}\, 8, 7, 2, 3)$ \eTD\eTR
\stopAlethe
\stopexample
\stopparagraph

\startparagraph{Skolemization and Other Preprocessing Steps.}
One typical example for a rule with context is the \proofRule{sko_ex}
rule that is used to express skolemization of an existentially
quantified variable.
%
The conclusion of a step that uses this rules is an equality.
The left-hand side is a formula starting with an existential quantifier
over some variable $x$.  In the formula on the right-hand side, the
variable is replaced by the appropriate Skolem term.
To provide a proof for the replacement, the \proofRule{sko_ex} step
uses one premise. 
The premise
has a context that maps the existentially quantified variables 
to the appropriate Skolem term.

\startAletheS
\setupTABLE[c][1][width=1em]
\setupTABLE[c][4][width=19em]
\setupTABLE[c][5][width=5em]
\bTR[bottomframe=on]\bTD[bottomframe=off] i. \eTD\bTD[leftframe=on] $\Gamma, x\mapsto (\varepsilon x.\,\varphi)$ \eTD\bTD \ctxsep
\eTD\bTD $\varphi ≈ \psi$ \eTD\bTD $(\dots)$ \eTD\eTR
\bTR\bTD j. \eTD\bTD $\Gamma$ \eTD\bTD \ctxsep
\eTD\bTD $(\exists x.\,\varphi) ≈ \psi$ \eTD\bTD $(\proofRule{sko_ex})$ \eTD\eTR
\stopAlethe

\startexample
To illustrate how such a rule is applied, consider the following example
taken from~\cite[barbosa-2019]. Here the term $\neg p(\varepsilon x.\neg
p(x))$ is skolemized. The \proofRule{refl} rule expresses
a simple tautology on the equality (reflexivity in this case), \proofRule{cong}
is functional congruence, and \proofRule{sko_forall} works like
\proofRule{sko_ex}, except that the choice term is $\varepsilon x.\neg\varphi$.

\startAletheS
\setupTABLE[c][1][width=1em]
\setupTABLE[c][4][width=17em]
\bTR\bTD 1. \eTD\bTD[leftframe=on] $x\mapsto (\varepsilon x.\,\neg (p\,x))$
     \eTD\bTD\ctxsep
     \eTD\bTD $x  ≈   \varepsilon x.\,\neg (p\,x) $ \eTD\bTD
     $\proofRule{refl}$ \eTD\eTR
\bTR[bottomframe=on]\bTD[bottomframe=off] 2. \eTD\bTD[leftframe=on]
     $x\mapsto (\varepsilon x.\,\neg (p\,x))$ \eTD\bTD\ctxsep \eTD\bTD $ (p\,x) ≈ p(\varepsilon x.\,\neg (p\,x))$ \eTD\bTD $(\proofRule{cong}\, 1)$ \eTD\eTR
\bTR\bTD 3. \eTD\bTD \eTD\bTD\ctxsep \eTD\bTD $(    \forall x.\,(p\,x))≈      (p\,(\varepsilon x.\,\neg (p\,x)))$ \eTD\bTD $(\proofRule{sko_forall}\, 2)$ \eTD\eTR
\bTR\bTD 4. \eTD\bTD \eTD\bTD\ctxsep \eTD\bTD $(\neg\forall x.\,(p\,x))≈ \neg (p\,(\varepsilon x.\,\neg (p\,x)))$ \eTD\bTD $(\proofRule{cong}\, 3)$ \eTD\eTR
\stopAlethe

\stopparagraph

\stopcomponent
